%!TeX root=MemoriaTFG.tex

\chapter*{Resumen}
\addcontentsline{toc}{chapter}{Resumen}
\markboth{RESUMEN}{RESUMEN}

¿Cuál es el rendimiento real y la capacidad de generalización de los
modelos fundacionales dermatológicos disponibles, en función del
volumen de datos etiquetados, del tipo de supervisión, de la
heterogeneidad taxonómica entre conjuntos de datos y de la
disparidad entre fototipos cutáneos? Esta pregunta articula el
presente trabajo. Su respuesta condiciona la viabilidad de integrar
estos modelos en flujos de teledermatología hospitalaria con
cohortes mediterráneas en español, escenario donde tres referencias
recientes liberadas por el grupo de la Universidad Monash
---PanDerm como codificador visual, Derm1M y la familia DermLIP como
modelos vision-lenguaje contrastivos, y DermFM-Zero como sistema
con capacidad zero-shot nativa--- no disponen de evaluación
independiente comparable a la fecha de cierre.

Para responderla, hemos evaluado catorce modelos fundacionales sobre
doce conjuntos de datos externos públicos, armonizados mediante una
ontología jerárquica de tres niveles ($4$ categorías etiológicas,
$43$ subcategorías clínicas y $367$ diagnósticos individuales)
codificada en SNOMED~CT y CIE-10. La evaluación combina cuatro
protocolos ---sondeo lineal, ajuste fino supervisado, segmentación
binaria con adaptación de bajo rango (LoRA) y clasificación
zero-shot multimodal--- ejecutados sobre una NVIDIA DGX Spark con
chip Grace Hopper GB10, $128$\,GB de memoria unificada y precisión
BF16, plataforma aarch64 no habitual en la literatura experimental.

PanDerm Large se confirma como el codificador congelado más fiable,
con una mejora media de $+9{,}0$\,pp en exactitud balanceada
respecto a PanDerm Base, y satura con un volumen reducido de datos
etiquetados: el $1\%$ del conjunto de entrenamiento de HAM10000
recupera el $96{,}2\%$ del AUROC obtenido con el conjunto completo.
La evaluación zero-shot multimodal queda $29{,}4$\,pp por debajo del
sondeo lineal, lo que cuantifica el coste de prescindir de
adaptación supervisada. La auditoría sistemática de solapamiento por
hash MD5 detecta fuga de información entre el preentrenamiento de
Derm1M y tres de los doce conjuntos de evaluación (Dermnet, HIBA y
MSKCC); cualquier cifra zero-shot reportada sobre ellos en la
literatura previa debe interpretarse acompañada de este
\emph{caveat}. La validación sobre el corpus DermapixelAI~1.0
---construido en paralelo a este trabajo--- confirma la
transferencia de los modelos a material clínico no anglosajón sin
reentrenamiento de la rama lingüística.

El trabajo aporta tres contribuciones reproducibles: un benchmark
comparativo homogéneo de las catorce familias bajo un mismo
protocolo, el dataset DermapixelAI~1.0 ---$1\,089$ imágenes
vinculadas a $669$ casos clínicos con texto narrativo y captions
sintéticas multi-perspectiva en castellano, distribuido bajo
licencia CC-BY-NC-SA~4.0--- y la auditoría sistemática de
solapamiento sobre los doce conjuntos de evaluación. Los tres
artefactos quedan disponibles para verificación independiente y
forman la base experimental sobre la que se construyen las líneas
de continuación documentadas en el capítulo de conclusiones.

\section*{Palabras clave}
modelos fundacionales, dermatología clínica, PanDerm, DermLIP,
evaluación empírica, ontología, equidad, heterogeneidad taxonómica.

% ----------------------------------------------------------------------

\chapter*{Abstract}
\addcontentsline{toc}{chapter}{Abstract}
\markboth{ABSTRACT}{ABSTRACT}

What is the actual performance and generalisation capacity of the
available dermatology foundation models, as a function of the
labelled data budget, the supervision regime, taxonomic
heterogeneity across datasets and skin phototype disparity? This
question articulates the present work. Its answer conditions the
viability of integrating these models into hospital teledermatology
workflows on Spanish-speaking Mediterranean cohorts, a scenario in
which three recent references released by the Monash group
---PanDerm as a vision encoder, Derm1M and the DermLIP family as
contrastive vision-language models, and DermFM-Zero as a system with
native zero-shot capability--- lack an independent comparable
evaluation at the time of writing.

To answer it, we evaluate fourteen foundation models over twelve
public external datasets, harmonised through a three-level
hierarchical ontology ($4$ aetiological families, $43$ clinical
subcategories and $367$ individual diagnoses) encoded in SNOMED~CT
and ICD-10. The evaluation combines four protocols ---linear
probing, supervised fine-tuning, binary lesion segmentation with
Low-Rank Adaptation (LoRA) and multimodal zero-shot
classification--- executed on a NVIDIA DGX Spark with a Grace
Hopper GB10 chip, $128$\,GB of unified memory and BF16 precision,
an aarch64 platform uncommon in the experimental literature.

PanDerm Large stands out as the most reliable frozen encoder, with a
mean improvement of $+9{,}0$\,pp in balanced accuracy over PanDerm
Base, and saturates with a small labelled budget: $1\%$ of the
HAM10000 training set recovers $96{,}2\%$ of the AUROC obtained with
the full set. Multimodal zero-shot evaluation falls $29{,}4$\,pp
below linear probing, bounding the cost of forgoing supervised
adaptation. The systematic MD5-based overlap audit detects
information leakage between the Derm1M pretraining set and three of
the twelve evaluation datasets (Dermnet, HIBA and MSKCC); any
zero-shot figure reported on them in prior literature must be
interpreted with this \emph{caveat}. Validation on the
DermapixelAI~1.0 corpus ---built in parallel with this work---
confirms model transfer to non-Anglophone clinical material without
retraining the language branch.

The work contributes three reproducible items: a homogeneous
comparative benchmark of the fourteen families under a common
protocol, the DermapixelAI~1.0 dataset ---$1\,089$ images linked to
$669$ clinical cases with narrative text and multi-perspective
synthetic captions in Spanish, released under a CC-BY-NC-SA~4.0
licence--- and the systematic overlap audit across the twelve
evaluation sets. The three artefacts remain available for
independent verification and form the experimental basis on which
the lines of continuation documented in the conclusions chapter are
built.

\section*{Keywords}
foundation models, clinical dermatology, PanDerm, DermLIP, empirical
evaluation, ontology, fairness, taxonomic heterogeneity.
